\chapter{SQL 解析实验}
\label{chap:parse}

\section{实验目的}
\label{sec:parse-purpose}

本章用于理解 DBMS\_C 如何把一条 SQL 字符串转换为其他模块可以使用的结构化对象。对于数据库系统来说，SQL 输入只是普通文本；在进入查询计划、扫描执行或更新执行之前，系统必须先完成词法切分、语法识别和 Data 对象构造。本章只讨论解析阶段，不执行 SQL，不生成 Plan，也不访问表数据。

本章继续采用“先补真实 CTest，再写章节”的实验组织方式。配套测试 \codepath{test/parse/ParseBasicTest.c} 覆盖 SELECT、INSERT、UPDATE、DELETE、CREATE TABLE、CREATE VIEW 和 CREATE INDEX 的基础解析行为。测试只断言 Parser 产出的 Data 对象是否符合输入 SQL，不进入第9章的 Plan/Scan，也不进入第10章的更新执行。

\section{SQL 解析在 DBMS\_C 中的位置}
\label{sec:parse-position}

SQL 解析位于命令行输入和执行模块之间。用户在 CLI 中输入 SQL 后，系统先通过 Lexer 将字符流切分为关键字、标识符、常量和分隔符，再由 Parser 按 SQL 类型构造对应的数据对象。例如 SELECT 会生成 \codepath{QueryData}，INSERT 会生成 \codepath{InsertData}，CREATE TABLE 会生成 \codepath{CreateTableData}。这些 Data 对象本身不负责执行，只负责保存“解析结果”。

图 \ref{fig:parse-sql-to-data} 展示了从 SQL 字符串到 Data 对象的转换关系。图中右侧的 Plan、UpdatePlanner 和 Transaction 是解析结果的下游消费者，本章不会展开它们的执行逻辑。

\begin{figure}[htbp]
  \centering
  \resizebox{0.94\textwidth}{!}{%
  \begin{tikzpicture}[node distance=8mm and 13mm]
    \node[dblevel] (sql) {SQL 字符串\\用户输入};
    \node[dbnode, right=of sql] (lexer) {Lexer\\读取 Token};
    \node[dbnode, right=of lexer] (parser) {Parser\\识别语法结构};
    \node[dbnode, below=of parser] (data) {Data 对象\\QueryData / CommandData};
    \node[dbsubnode, below left=of data] (plan) {第9章\\查询计划与扫描};
    \node[dbsubnode, below=of data] (update) {第10章\\更新执行};
    \node[dbsubnode, below right=of data] (meta) {第7章\\元数据支持};

    \draw[dbarrow] (sql) -- (lexer);
    \draw[dbarrow] (lexer) -- (parser);
    \draw[dbdown] (parser) -- (data);
    \draw[dbdown] (data) -- (plan);
    \draw[dbdown] (data) -- (update);
    \draw[dbdown] (data) -- (meta);
  \end{tikzpicture}}
  \caption{SQL 字符串到 Data 对象的转换关系}
  \label{fig:parse-sql-to-data}
\end{figure}

\section{相关源码文件}
\label{sec:parse-source-files}

表 \ref{tab:parse-source-files} 列出本章直接涉及的真实源码文件。阅读时建议先看 \codepath{Lexer.h} 和 \codepath{Parser.h}，确认对外 API；再看 \codepath{Parser.c}，理解不同 SQL 命令如何被分发；最后结合各 Data 对象头文件观察解析结果保存在哪里。

\begin{table}[htbp]
  \centering
  \caption{Parse 模块相关源码文件}
  \label{tab:parse-source-files}
  \begin{tabularx}{\textwidth}{l Y}
    \toprule
    文件 & 本章关注点 \\
    \midrule
    \codepath{parse/Lexer.h}, \codepath{parse/Lexer.c} & 基于 \codepath{StreamTokenizer} 读取关键字、标识符、整数常量、字符串常量和分隔符。\\
    \codepath{parse/Parser.h}, \codepath{parse/Parser.c} & SQL 解析主入口，构造 \codepath{QueryData} 或 \codepath{CommandData}。\\
    \codepath{parse/PredParser.h}, \codepath{parse/PredParser.c} & 谓词解析相关结构；当前主要逻辑在 Parser 的 \apiname{ParserPredicate}、\apiname{ParserTerm} 中体现。\\
    \codepath{parse/QueryData.h}, \codepath{parse/QueryData.c} & 保存 SELECT 字段列表、表名列表和谓词。\\
    \codepath{parse/InsertData.h}, \codepath{parse/InsertData.c} & 保存 INSERT 的表名、字段列表和值列表。\\
    \codepath{parse/ModifyData.h}, \codepath{parse/ModifyData.c} & 保存 UPDATE 的表名、目标字段、新值表达式和谓词。\\
    \codepath{parse/DeleteData.h}, \codepath{parse/DeleteData.c} & 保存 DELETE 的表名和谓词。\\
    \codepath{parse/CreateTableData.h}, \codepath{parse/CreateTableData.c} & 保存 CREATE TABLE 的表名和 Schema。\\
    \codepath{parse/CreateViewData.h}, \codepath{parse/CreateViewData.c} & 保存 CREATE VIEW 的视图名和内部 SELECT 查询定义。\\
    \codepath{parse/CreateIndexData.h}, \codepath{parse/CreateIndexData.c} & 保存 CREATE INDEX 的索引名、表名和字段名。\\
    \codepath{query/Predicate.h}, \codepath{query/Term.h}, \codepath{query/Expression.h} & 表示 WHERE 条件中的谓词、比较项和表达式。\\
    \bottomrule
  \end{tabularx}
\end{table}

\section{Lexer、Parser 与 PredParser 的职责}
\label{sec:parse-lexer-parser}

Lexer 的职责是面向 Token：它不理解整条 SQL 的业务含义，只负责判断当前 Token 是否匹配某类符号，并在需要时消费它。Parser 的职责是面向语法结构：它按照 SQL 命令格式调用 Lexer，逐步构造字段列表、表名列表、常量列表、表达式、谓词和命令 Data 对象。PredParser 当前只保留了谓词解析结构，实际谓词构造主要通过 Parser 中的 \apiname{ParserPredicate} 和 \apiname{ParserTerm} 完成。

图 \ref{fig:parse-lexer-parser-pred} 展示三者的关系。注意，本章只观察 Parser 产物，不尝试让 Parser 执行 SQL。

\begin{figure}[htbp]
  \centering
  \resizebox{0.88\textwidth}{!}{%
  \begin{tikzpicture}[node distance=8mm and 13mm]
    \node[dblevel] (st) {StreamTokenizer\\底层词法切分};
    \node[dbnode, right=of st] (lexer) {Lexer\\Match / Eat Token};
    \node[dbnode, right=of lexer] (parser) {Parser\\SQL 语法组合};
    \node[dbnode, below=of parser] (pred) {Predicate / Term\\WHERE 条件对象};
    \node[dbsubnode, below left=of pred] (expr) {Expression\\字段名或常量};
    \node[dbsubnode, below right=of pred] (constant) {Constant\\int 或 string};

    \draw[dbarrow] (st) -- (lexer);
    \draw[dbarrow] (lexer) -- (parser);
    \draw[dbdown] (parser) -- (pred);
    \draw[dbdown] (pred) -- (expr);
    \draw[dbdown] (pred) -- (constant);
  \end{tikzpicture}}
  \caption{Lexer、Parser 与谓词对象的关系}
  \label{fig:parse-lexer-parser-pred}
\end{figure}

\section{不同 SQL 命令的解析结果}
\label{sec:parse-data-objects}

Parser 对 SELECT 和更新类命令使用不同入口。SELECT 通过 \apiname{ParserQuery} 返回 \codepath{QueryData}；INSERT、DELETE、UPDATE 和 CREATE 族命令通过 \apiname{ParserUpdateCmd} 返回 \codepath{CommandData}。表 \ref{tab:parse-command-data-map} 汇总了本章测试覆盖的 SQL 与 Data 对象映射。

\begin{table}[htbp]
  \centering
  \caption{SQL 命令到 Data 对象的映射}
  \label{tab:parse-command-data-map}
  \begin{tabularx}{\textwidth}{l l Y}
    \toprule
    SQL 类型 & 解析结果 & 保存的核心信息 \\
    \midrule
    SELECT & \codepath{QueryData} & \codepath{fields}、\codepath{tables}、\codepath{predicate}\\
    INSERT & \codepath{CommandData} + \codepath{InsertData} & 表名、字段列表、常量值列表\\
    UPDATE & \codepath{CommandData} + \codepath{ModifyData} & 表名、目标字段、新值表达式、谓词\\
    DELETE & \codepath{CommandData} + \codepath{DeleteData} & 表名、谓词\\
    CREATE TABLE & \codepath{CommandData} + \codepath{CreateTableData} & 表名、Schema 字段类型和长度\\
    CREATE VIEW & \codepath{CommandData} + \codepath{CreateViewData} & 视图名、内部 SELECT 的 QueryData\\
    CREATE INDEX & \codepath{CommandData} + \codepath{CreateIndexData} & 索引名、表名、字段名\\
    \bottomrule
  \end{tabularx}
\end{table}

\section{配套测试}
\label{sec:parse-test}

本章新增并注册了配套测试文件：

\begin{itemize}
  \item \codepath{test/parse/ParseBasicTest.c}
\end{itemize}

该测试不创建数据库目录，也不访问文件、记录、元数据或执行器。它只构造 Parser，输入短 SQL 字符串，然后断言返回的结构体字段符合预期。表 \ref{tab:parse-tests} 列出每个测试函数的覆盖内容。

\begin{table}[htbp]
  \centering
  \caption{ParseBasicTest 覆盖内容}
  \label{tab:parse-tests}
  \begin{tabularx}{\textwidth}{l Y}
    \toprule
    测试函数 & 主要断言 \\
    \midrule
    \codepath{test_parse_select_fields_tables_and_predicate} & \codepath{QueryData} 非空；字段列表为 \codepath{id,name}；表名为 \codepath{student}；谓词为 \codepath{id = 1}。\\
    \codepath{test_parse_insert_table_fields_and_values} & \codepath{CommandData.code} 为 \codepath{CMD_INSERT_DATA}；表名、字段数量、值数量、整数字面量和字符串字面量正确。\\
    \codepath{test_parse_update_target_value_and_predicate} & \codepath{CMD_MODIFY_DATA}；表名为 \codepath{student}；目标字段为 \codepath{name}；新值为字符串常量；谓词为 \codepath{id = 1}。\\
    \codepath{test_parse_delete_table_and_predicate} & \codepath{CMD_DELETE_DATA}；表名为 \codepath{student}；谓词对象非空且表示 \codepath{id = 1}。\\
    \codepath{test_parse_create_table_schema} & \codepath{CMD_CREATE_TABLE}；表名为 \codepath{student}；Schema 含 \codepath{id int} 和 \codepath{name varchar(20)}。\\
    \codepath{test_parse_create_view_query_data} & \codepath{CMD_CREATE_VIEW}；视图名正确；内部 \codepath{QueryData} 的字段、表名和谓词正确。\\
    \codepath{test_parse_create_index_identifiers} & \codepath{CMD_CREATE_INDEX}；索引名、表名和字段名正确。\\
    \bottomrule
  \end{tabularx}
\end{table}

本章测试已经在临时验证目录 \codepath{build-check} 中通过。单独运行新增测试的命令如下：

\begin{dbcode}[listing options={language=bash}]{运行 ParseBasicTest}
ctest --test-dir build-check --output-on-failure -R ParseBasicTest
\end{dbcode}

运行 Parse 相关测试的命令如下：

\begin{dbcode}[listing options={language=bash}]{运行 Parse 相关测试}
ctest --test-dir build-check --output-on-failure -R Parse
\end{dbcode}

\section{关键代码讲解}
\label{sec:parse-code}

第一段代码来自 \codepath{parse/Parser.c} 的 \apiname{ParserQuery}。它说明 SELECT 语句如何被拆成字段列表、表名列表和可选 WHERE 谓词。

\begin{dbcode}{ParserQuery 构造 QueryData}
QueryData *ParserQuery(Parser*parser){
    LexerEatKeyWord(parser->lexer,"select");
    CList *fields = ParserSelectList(parser);
    LexerEatKeyWord(parser->lexer,"from");
    CList*tables = ParserTabletList(parser);
    Predicate *predicate = PredicateInit(NULL);
    if(LexerMatchKeyWord(parser->lexer,"where")){
        LexerEatKeyWord(parser->lexer,"where");
        predicate = ParserPredicate(parser);
    }
    return QueryDataInit(fields,tables,predicate);
}
\end{dbcode}

第二段代码来自 \codepath{parse/Parser.c} 的 \apiname{ParserUpdateCmd}。它是更新类和 CREATE 类命令的统一分发入口，根据当前关键字选择具体解析函数。

\begin{dbcode}{ParserUpdateCmd 分发更新类命令}
CommandData *ParserUpdateCmd(Parser*parser){
    if(LexerMatchKeyWord(parser->lexer,"insert")){
        return ParserInsert(parser);
    }else if(LexerMatchKeyWord(parser->lexer,"delete")){
        return ParserDelete(parser);
    }else if(LexerMatchKeyWord(parser->lexer,"update")){
        return ParserModify(parser);
    }else{
        return ParserCreate(parser);
    }
}
\end{dbcode}

第三段代码来自 \codepath{parse/Parser.c} 的 \apiname{ParserTerm} 和 \apiname{ParserPredicate}。当前实现把一个比较表达式封装为 Term；如果遇到 \codepath{and}，再把后面的 Predicate 合并进来。

\begin{dbcode}{ParserTerm 与 ParserPredicate 构造 WHERE 条件}
Term *ParserTerm(Parser*parser){
    Expression *lhs = ParserExpression(parser);
    const CompareOp op = LexerGetCurrentOp(parser->lexer);
    LexerNextToken(parser->lexer);
    Expression *rhs = ParserExpression(parser);
    return TermInit(lhs,rhs,op);
}

Predicate *ParserPredicate(Parser*parser){
    Predicate * predicate = PredicateInit(ParserTerm(parser));
    if(LexerMatchKeyWord(parser->lexer,"and")){
        LexerEatKeyWord(parser->lexer,"and");
        PredicateConjoinWith(predicate, ParserPredicate(parser));
    }
    return predicate;
}
\end{dbcode}

第四段代码来自 \codepath{parse/Parser.c} 的 \apiname{ParserCreateTable}。它把字段定义列表解析为 \codepath{Schema}，再包装进 \codepath{CreateTableData}。

\begin{dbcode}{ParserCreateTable 构造 CreateTableData}
CommandData* ParserCreateTable(Parser*parser){
    LexerEatKeyWord(parser->lexer,"table");
    CString *tblname = CStringCreateFromCStr(LexerEatId(parser->lexer));
    LexerEatDelim(parser->lexer,'(');
    Schema *schema = ParserFieldDefs(parser);
    LexerEatDelim(parser->lexer,')');

    CommandData * commandData = malloc(sizeof(CommandData));
    commandData->code = CMD_CREATE_TABLE;
    commandData->data.createTableData = CreateTableDataInit(tblname,schema);
    return commandData;
}
\end{dbcode}

\section{运行与观察}
\label{sec:parse-run}

本章观察重点不是 SQL 是否能产生结果，而是解析出的结构体是否正确。运行 \codepath{ParseBasicTest} 时，可以重点观察以下几个方面：

\begin{enumerate}
  \item SELECT 是否被解析为 \codepath{QueryData}，并保存字段列表、表名列表和谓词对象。
  \item INSERT、UPDATE、DELETE、CREATE 族命令是否被统一包装为 \codepath{CommandData}，并设置正确的 \codepath{CommandCode}。
  \item 字符串常量和整数常量是否被区分为不同的 \codepath{Constant}。
  \item \codepath{id = 1} 这样的 WHERE 条件是否被表示为 \codepath{Predicate -> Term -> Expression/Constant}。
  \item CREATE TABLE 是否生成真实 \codepath{Schema}，并能读出字段类型和长度。
\end{enumerate}

图 \ref{fig:parse-command-flow} 用 INSERT 和 CREATE TABLE 为例说明解析阶段的调用链。它只展示 Data 对象生成，不展示执行阶段。

\begin{figure}[htbp]
  \centering
  \resizebox{0.96\textwidth}{!}{%
  \begin{tikzpicture}[node distance=8mm and 11mm]
    \node[dblevel] (cmd) {ParserUpdateCmd};
    \node[dbnode, below left=of cmd] (insert) {ParserInsert\\INSERT};
    \node[dbnode, below=of cmd] (modify) {ParserModify / ParserDelete\\UPDATE / DELETE};
    \node[dbnode, below right=of cmd] (create) {ParserCreate\\CREATE};
    \node[dbsubnode, below=of insert] (insertdata) {InsertData\\tblname, flds, vals};
    \node[dbsubnode, below=of modify] (moddata) {ModifyData / DeleteData\\predicate};
    \node[dbsubnode, below=of create] (createdata) {CreateTable/View/IndexData};

    \draw[dbdown] (cmd) -- (insert);
    \draw[dbdown] (cmd) -- (modify);
    \draw[dbdown] (cmd) -- (create);
    \draw[dbdown] (insert) -- (insertdata);
    \draw[dbdown] (modify) -- (moddata);
    \draw[dbdown] (create) -- (createdata);
  \end{tikzpicture}}
  \caption{更新类 SQL 命令到 Data 对象的解析调用链}
  \label{fig:parse-command-flow}
\end{figure}

\section{思考题}
\label{sec:parse-questions}

\begin{enumerate}
  \item 为什么 Parser 不直接执行 SQL，而是先生成 \codepath{QueryData} 或 \codepath{CommandData}？
  \item \codepath{LexerMatchKeyWord} 和 \codepath{LexerEatKeyWord} 的区别是什么？如果把二者混用，会产生什么问题？
  \item WHERE 条件为什么要拆成 \codepath{Predicate}、\codepath{Term} 和 \codepath{Expression} 三层对象？
  \item CREATE TABLE 解析阶段为什么只生成 \codepath{Schema}，而不直接创建表文件？
\end{enumerate}
